% Confirma o intervalo da caixa de foliação manual.
\documentclass[permitir-placeholders]{sindicancia}
% ============================================================
% DADOS GERAIS DA SINDICÂNCIA
% Arquivo: dados/dados-sindicancia.tex
%
% Este arquivo deve existir sempre.
% Ele concentra SOMENTE os dados gerais reutilizados por mais
% de um modelo/anexo.
%
% Regra do projeto:
% - Dados comuns ficam aqui.
% - Dados exclusivos de uma peça ficam em dados-anexo-*.tex.
% - Evite redefinir estas macros nos arquivos de anexos, pois
%   o LaTeX usa sempre o último valor carregado no main.tex.
% ============================================================

% ------------------------------------------------------------
% Data/local de abertura dos trabalhos
% Usada também como fallback quando um documento não informar data própria.
% ------------------------------------------------------------
\dia{00}
\mes{mês}
\ano{2026}
\cidade{CIDADE/UF}

% ------------------------------------------------------------
% Organização militar
% ------------------------------------------------------------
\om{ORGANIZAÇÃO MILITAR}
\omSuperior{ESCALÃO SUPERIOR}
\omComplemento{LINHA COMPLEMENTAR DA OM}
\omNomeHistorico{NOME HISTÓRICO DA OM}

% ------------------------------------------------------------
% Processo / portaria / objeto
% ------------------------------------------------------------
\nup{00000.000000/2026-00}

\portaria{Portaria nº 000, de 00 de mês de 2026, do Comandante da ORGANIZAÇÃO MILITAR}

\objeto{apurar os fatos constantes da portaria instauradora e seus anexos}

% ------------------------------------------------------------
% Autoridade instauradora / nomeante
% ------------------------------------------------------------
\autoridadeNome{AUTORIDADE INSTAURADORA}
\autoridadePosto{POSTO}
\autoridadeFuncao{Comandante da ORGANIZAÇÃO MILITAR}

% ------------------------------------------------------------
% Sindicante
% ------------------------------------------------------------
\sindicanteNome{SINDICANTE}
\sindicantePosto{POSTO/GRADUAÇÃO}
\sindicanteOM{\getOm}

% ------------------------------------------------------------
% Sindicado
% Se ainda não houver sindicado, deixe em branco ou comente,
% desde que o modelo usado permita ausência.
% ------------------------------------------------------------
\sindicadoNome{SINDICADO}
\sindicadoPosto{POSTO/GRADUAÇÃO}
\sindicadoSU{SUBUNIDADE/SEÇÃO}
\sindicadoOM{\getOm}

% Dados complementares do sindicado, usados principalmente no Anexo Q
% e também por outros modelos que precisem qualificar o sindicado.
\sindicadoIdentidade{000000000 ÓRGÃO/UF}
\sindicadoCPF{000.000.000-00}
\sindicadoNascimento{00 de mês de 0000}
\sindicadoNaturalidade{CIDADE/UF}
\sindicadoEstadoCivil{ESTADO CIVIL}
\sindicadoFiliacao{NOME DO PAI e NOME DA MÃE}
\sindicadoResidencia{ENDEREÇO COMPLETO}
\sindicadoCelular{(00) 00000-0000}

% ------------------------------------------------------------
% Escrivão
% Usado pelos Anexos F e G, quando houver designação de escrivão.
% ------------------------------------------------------------
\escrivaoNome{ESCRIVÃO}
\escrivaoPosto{POSTO/GRADUAÇÃO}

% ------------------------------------------------------------
% Campos genéricos opcionais
% ------------------------------------------------------------
\prazo{3 (três) dias úteis}


\makeatletter
\newcommand{\ExigirCaixaOculta}[1]{%
  \ifsind@mostrarflbox
    \ClassError{sindicancia}{Caixa de folha visivel em #1}{A caixa deveria estar oculta neste ponto.}%
  \fi
}
\newcommand{\ExigirCaixaVisivel}[1]{%
  \ifsind@mostrarflbox\else
    \ClassError{sindicancia}{Caixa de folha oculta em #1}{A caixa deveria estar visivel neste ponto.}%
  \fi
}
\makeatother

\begin{document}
\ExigirCaixaOculta{estado inicial}

% DADOS - ANEXO A
\dataDocumento{1}{maio}{2026}{CIDADE/UF}
\chaveCronologica{20260501}
\modoDocumento{padrao}
\portariaNumero{001/2026}
\portariaAnexoOrigem{DIEx nº 001, de 30 de abril de 2026}
\portariaSinteseFatos{fato de interesse da Administração Militar descrito no documento anexo}
\portariaPrazoDias{30}
% Opcional, apenas quando a complexidade recomendar:
% \portariaEscrivao{POSTO/GRADUAÇÃO NOME DO ESCRIVÃO}
\portariaPublicacaoBI{BI nº 080, de 1º de maio de 2026}

\GerarDocumento{portaria-instauracao}
\ExigirCaixaOculta{anexo A}

% ============================================================
% DADOS - ANEXO C - CAPA DOS AUTOS
% Normalmente estes campos ficam em dados/dados-sindicancia.tex.
% ============================================================

\nup{64108.000000/2026-00}
\objeto{apurar os fatos constantes da documentação inicial da sindicância.}

\sindicanteNome{RAFAEL DOS SANTOS DANTAS}
\sindicantePosto{1º Ten}

% Preencha apenas se houver sindicado definido.
\sindicadoNome{NOME DO SINDICADO}
\sindicadoPosto{POSTO/GRADUAÇÃO}

\GerarDocumento{capa}
\ExigirCaixaOculta{capa}

% ============================================================
% DADOS - ANEXO D - TERMO DE ABERTURA
% Arquivo: dados/dados-anexo-d-termo-abertura.tex
%
% Este anexo usa apenas dados gerais definidos em:
%   dados/dados-sindicancia.tex
%
% Não redefina aqui macros como:
%   \dia
%   \mes
%   \ano
%   \cidade
%   \om
%   \portaria
%   \sindicanteNome
%   \sindicantePosto
%
% O arquivo é mantido para preservar a organização por anexo
% e permitir comentários específicos do termo de abertura, se necessário.
% ============================================================

\GerarDocumento{termo-abertura}
\ExigirCaixaVisivel{termo de abertura}

% ============================================================
% DADOS - ANEXO H - DESPACHO 1
% Salvar como dados/dados-despacho1.tex para uso real.
% ============================================================

\novoDespacho

\despachoNumero{1}

\despachoTexto{%
Oficiar ao Sr. Comandante do(a) ORGANIZAÇÃO MILITAR, solicitando a remessa de cópia de DOCUMENTO/INFORMAÇÃO necessária à presente sindicância.
}

% Se quiser data própria:
% \despachoLocalData{Petrolina/PE, 31 de maio de 2026}

\GerarDocumento{despacho}
\ExigirCaixaVisivel{instrucao}

% ============================================================
% DADOS - ANEXO Y - DIEx DE REMESSA
% Arquivo: dados/dados-anexo-y-diex-remessa.tex
% ============================================================

% ------------------------------------------------------------
% Data própria do documento
% Deixe em branco para usar a data/local do termo de abertura.
% Formato: \dataDocumento{dia}{mês}{ano}{cidade/UF}
% ------------------------------------------------------------
\dataDocumento{}{}{}{}

\novoDiexRemessa

\diexRemessaNumero{003}
\diexRemessaOrigem{Sindicante}

% Se não informar, usa \localData.
% \diexRemessaLocalData{CIDADE/UF, 31 de maio de 2026}

\diexRemessaDestinatario{\getAutoridadePosto\ \getAutoridadeNome, \getAutoridadeFuncao}
% A quantidade de folhas no assunto do DIEx de remessa é deixada em branco
% no PDF para preenchimento manual no fechamento dos autos.

% Se não informar, usa \getPortaria.
% \diexRemessaReferencia{Portaria nº 000, de 00 de mês de 2026}

% Opcional:
% \diexRemessaTextoComplementar{Informo que os autos seguem em anexo para apreciação e solução por essa autoridade instauradora.}

\GerarDocumento{diex-remessa}
\ExigirCaixaOculta{DIEx de remessa}

% DADOS - ANEXO Z
% Referência da Seção de Apoio Jurídico; não é elaborado pelo sindicante.
\dataDocumento{25}{maio}{2026}{CIDADE/UF}
\chaveCronologica{20260525}
\modoDocumento{padrao}
\solucaoTipo{acolher}
\solucaoParecer{os fatos foram esclarecidos nos termos do relatório}
\solucaoFundamentos{%
  \begin{enumerate}[label=\alph*)]
    \item a instrução observou as formalidades aplicáveis;
    \item as provas foram analisadas de forma comparativa e motivada;
    \item foi assegurado o contraditório quando presente a figura do sindicado.
  \end{enumerate}
}
\solucaoProvidencias{%
  \begin{enumerate}[label=\alph*)]
    \item arquivamento dos autos;
    \item publicação em boletim interno;
    \item ciência aos interessados cabíveis.
  \end{enumerate}
}
\solucaoPublicacaoBI{BI nº 096, de 25 de maio de 2026}
\solucaoNotificados{sindicado e denunciante/ofendido, quando houver}

\GerarDocumento{solucao-sindicancia}
\ExigirCaixaOculta{anexo Z}
\end{document}
